\chapter{Hydronephrosis}
\index{Hydronephrosis}
\label{sec:Hydronephrosis}

\section{Overview}
Hydronephrosis is the dilation of the renal pelvis and calyces caused by obstruction to urine outflow from the kidney. It can be unilateral or bilateral and may result from intrinsic obstruction, extrinsic compression, or functional disorders of the urinary tract. If left untreated, hydronephrosis can lead to progressive renal parenchymal atrophy and loss of kidney function.

\section{Classification}
\begin{description}[font=\normalfont\bfseries,leftmargin=3em,labelwidth=2.5em]
  \item[Cause] Obstruction of urine flow at any level from the renal pelvis to the urethral meatus
  \item[Organ System] Kidneys
  \item[Pathophysiology] Obstruction increases intraluminal pressure proximal to the blockage, which is transmitted to the renal tubules. This causes tubular dilation, interstitial inflammation, reduced glomerular filtration rate, and progressive nephron loss. Chronic obstruction leads to renal parenchymal thinning and irreversible fibrosis.
  \item[Duration] Acute or chronic
  \item[Onset] Sudden (acute obstruction) or gradual (chronic partial obstruction)
  \item[Mode of Transmission] Not applicable
  \item[Clinical Course] Depends on the degree and duration of obstruction; may be asymptomatic or present with flank pain, urinary symptoms, and declining renal function.
  \item[Severity] Mild to severe; can be life-threatening if bilateral and complete
  \item[Extent of Disease] Unilateral or bilateral
  \item[Age Group] All ages; etiology varies by age (congenital in children, acquired in adults)
  \item[Epidemiology] Common finding on prenatal ultrasound (1--5\% of pregnancies). In adults, prevalence increases with age due to benign prostatic hyperplasia, nephrolithiasis, and malignancy.
  \item[ICD-11 Code] GB00.0
\end{description}

\section{Causes}
Hydronephrosis results from obstruction at any level of the urinary tract. Common causes include ureteropelvic junction obstruction, ureteral stones (nephrolithiasis), benign prostatic hyperplasia, prostate cancer, retroperitoneal fibrosis, ureteral strictures, and neurogenic bladder. In children, congenital anomalies such as posterior urethral valves or ureteropelvic junction obstruction are common.

\section{Risk Factors}
\begin{itemize}[noitemsep]
  \item Nephrolithiasis (kidney stones)
  \item Benign prostatic hyperplasia (older men)
  \item Pelvic or abdominal malignancies (cervical, prostate, bladder, colorectal)
  \item Congenital urinary tract anomalies
  \item Previous pelvic surgery or radiation therapy
  \item Spinal cord injury or neurogenic bladder
\end{itemize}

\section{Signs and Symptoms}
\begin{itemize}[noitemsep]
  \item Flank pain (dull ache or acute colic with stone obstruction)
  \item Nausea and vomiting
  \item Decreased urine output (oliguria or anuria) in bilateral obstruction
  \item Hematuria (if caused by stones or infection)
  \item Urinary tract infections or pyelonephritis
  \item Difficulty performing daily activities
\end{itemize}

\section{Progression and Stages}
Acute complete obstruction causes rapid rise in intrarenal pressure, severe pain, and potentially irreversible damage within hours to days. Chronic partial obstruction leads to gradual calyceal dilation, progressive cortical thinning, and declining renal function over months to years. The grade of hydronephrosis (mild, moderate, severe) correlates with the degree of pelvicalyceal dilation and parenchymal loss.

\section{Complications}
\begin{itemize}[noitemsep]
  \item Renal parenchymal atrophy and irreversible kidney damage
  \item Chronic kidney disease or end-stage renal disease (bilateral obstruction)
  \item Urinary tract infection and pyelonephritis
  \item Hypertension (from renin-angiotensin activation)
  \item Post-obstructive diuresis after relief of obstruction
  \item Reduced quality of life
  \item Emotional and psychological effects
\end{itemize}

\section{Diagnosis}
\begin{itemize}[noitemsep]
  \item Renal ultrasound (initial imaging of choice to assess hydronephrosis grade)
  \item CT urography or MRI urography for detailed anatomical evaluation
  \item Lab tests (serum creatinine, BUN, electrolytes, urinalysis)
  \item Imaging tests (diuretic renography to quantify obstruction, retrograde pyelography)
\end{itemize}

\section{Prevention}
\begin{itemize}[noitemsep]
  \item Maintain adequate hydration to reduce stone formation
  \item Prompt treatment of urinary tract infections
  \item Screening for congenital anomalies in at-risk pregnancies
  \item Regular check-ups
\end{itemize}

\section{Treatment Options}
\begin{itemize}[noitemsep]
  \item Relief of obstruction (urethral catheter, ureteral stent, nephrostomy tube)
  \item Treatment of underlying cause (extracorporeal shockwave lithotripsy for stones, transurethral resection of prostate for BPH)
  \item Surgical correction (pyeloplasty for UPJ obstruction, ureteral reimplantation)
  \item Antibiotics for associated urinary tract infection
  \item Lifestyle changes
\end{itemize}

\section{Living with It}
\begin{itemize}[noitemsep]
  \item Follow treatment plan
  \item Regular appointments
  \item Adjust daily habits
  \item Support groups
  \item Keep learning
\end{itemize}

\section{Outlook}
Outcome depends on the severity, duration, and etiology of the obstruction. Acute unilateral hydronephrosis relieved promptly usually has an excellent prognosis with full recovery of renal function. Chronic or high-grade obstruction may cause irreversible renal damage, with 10--30\% of patients developing long-term renal impairment. Bilateral complete obstruction is a medical emergency requiring urgent decompression to prevent permanent renal failure.
